%----------------------------------------------------------------------------------------
%	RESEARCH
%----------------------------------------------------------------------------------------

\section{Research Experience}

\researchentry
  {Starspot Inference with Light Curve Inversion Techniques}
  {American Museum of Natural History}
  {M.S. Thesis Researcher}
  {Dr. Lucy Lu}
  {Aug 2023 -- Sep 2025}
  {\thesislink{https://academicworks.cuny.edu/gc_etds/6480/}}
\begin{itemize}[
  leftmargin=1em,
  label={--},
  itemsep=3pt,
  topsep=0pt,
  parsep=0pt,
  partopsep=0pt
]
  \item Implemented forward- and inverse-modeling workflows with
  \texttt{fleck} and \texttt{starry} to analyze 1,000 synthetic, evolving
  starspot light curves generated with \texttt{butterpy}.

  \item Evaluated recovery of stellar rotation, activity indicators, and
  surface structure across variations in photometric noise, viewing
  inclination, and spot configuration.

  \item Recovered rotation and statistical activity indicators in the tested
  light curves, while finding that individual starspot locations and surface
  maps were not uniquely constrained.
\end{itemize}

\vspace{0.4em}

\researchentry
  {Towards Mapping Brown Dwarf and Giant Exoplanet Atmospheres}
  {American Museum of Natural History}
  {BDNYC Researcher}
  {Dr. Johanna Vos}
  {May 2021 -- Aug 2023}
  {\githubrepo{https://github.com/MohammadRefat23/StarryAtmospheres}}

\begin{itemize}[
  leftmargin=1em,
  label={--},
  itemsep=3pt,
  topsep=0pt,
  parsep=0pt,
  partopsep=0pt
]

\item Applied spherical-harmonic surface mapping with \texttt{starry} to
three-dimensional atmospheric simulations and 20-hour Spitzer light curves
of two isolated planetary-mass objects.

\item Evaluated model complexity, identifying spherical-harmonic degree
$l=5$ as sufficient to reproduce the relevant light-curve structure.

\item Identified atmospheric evolution as a limitation of static-map fits,
motivating rotation-by-rotation mapping of changing cloud structures.

\end{itemize}

\vspace{0.4em}

\researchentry
  {Chemodynamically Characterizing the Jhelum Stellar Stream with APOGEE-2}
  {Sloan Digital Sky Survey Faculty and Student Team}
  {SDSS FAST Student Researcher}
  {Dr. Allyson Sheffield}
  {May 2020 -- May 2021}
  {%
    \githubrepo{https://github.com/MohammadRefat23/jhelum-stellar-stream}
    \hspace{0.15em}
    \paperlink{https://doi.org/10.3847/1538-4357/abee93}
  }

\begin{itemize}[
  leftmargin=1em,
  label={--},
  itemsep=3pt,
  topsep=0pt,
  parsep=0pt,
  partopsep=0pt
]
\item Combined APOGEE-2 spectroscopy, Gaia astrometry, stellar kinematics,
and chemical abundances to identify candidate Jhelum stream members in a
sample of 289 red giants.

\item Contributed to identifying a new metal-poor red-giant member and
compared its chemical and orbital properties with candidate
Gaia--Enceladus--Sausage stars.

\item Coauthored a peer-reviewed article in \textit{The Astrophysical Journal}
on the stream's membership, chemical properties, and progenitor origin.
\end{itemize}

\vspace{0.4em}

\researchentry
  {Looking at Star Formation Through Chemistry}
  {American Museum of Natural History}
  {AstroCom NYC Scholar}
  {Dr. Allyson Sheffield}
  {Aug 2018 -- May 2020}
  {}

\begin{itemize}[
  leftmargin=1em,
  label={--},
  itemsep=3pt,
  topsep=0pt,
  parsep=0pt,
  partopsep=0pt
]
\item Processed and normalized high-resolution echelle spectra with
\texttt{SMHr}, measured absorption-line equivalent widths, and iteratively
fitted stellar atmospheric parameters and elemental abundances for
chemical-tagging analysis.

\item Validated the workflow against published Arcturus measurements,
recovering effective temperature within 50 K and metallicity within
0.07 dex of the reference values.
\end{itemize}

\vspace{0.4em}

\researchentry
  {Two-Point Statistics in the Star-Forming Interstellar Medium}
  {Flatiron Institute, Center for Computational Astrophysics}
  {AstroCom NYC Scholar}
  {Dr. Chang-Goo Kim}
  {Jun 2018 -- Aug 2018}
  {\githubrepo{https://github.com/MohammadRefat23/tigress-metallicity-correlations}}

\begin{itemize}[
  leftmargin=1em,
  label={--},
  itemsep=3pt,
  topsep=0pt,
  parsep=0pt,
  partopsep=0pt
]
\item Developed Python routines to compute and visualize two-point
metallicity correlations in TIGRESS magnetohydrodynamic simulations of the
star-forming interstellar medium.

\item Measured a characteristic metallicity half-correlation length of
$120\pm25\,\mathrm{pc}$ and estimated a turbulent metal-mixing timescale of
$\sim10\,\mathrm{Myr}$.
\end{itemize}

\researchentry
  {Staying in Science}
  {American Museum of Natural History}
  {Research Contributor, NSF-Funded Study}
  {Dr. Preeti Gupta}
  {Aug 2017 -- Sep 2019}
  {}

\begin{itemize}[
  leftmargin=1em,
  label={--},
  itemsep=3pt,
  topsep=0pt,
  parsep=0pt,
  partopsep=0pt
]
\item Prepared research memos and contributed participant information used
to develop social-network maps of students' STEM education and career
pathways.

\item Reviewed survey data to examine STEM persistence among students from
historically marginalized backgrounds who participated in science research
mentoring programs.
\end{itemize}

\vspace{0.4em}